\documentclass[aspectratio=1610]{beamer}

\usetheme{default}
\usecolortheme{seahorse}
\beamertemplatenavigationsymbolsempty
\setbeamertemplate{footline}[frame number]

% \definecolor{alert}{HTML}{B16286}
% \definecolor{highlight}{HTML}{458588}
% \definecolor{gruvbox_light}{HTML}{F9F5D7}
% \definecolor{gruvbox_dark}{HTML}{3C3836}

% \usecolortheme[named=highlight]{structure}
% \setbeamercolor{normal text}{fg=gruvbox_dark}
% \setbeamercolor{alerted text}{fg=alert}

\beamertemplatenavigationsymbolsempty
\setbeamertemplate{footline}[frame number]

\setbeamercovered{transparent}

\titlegraphic{
  \begin{center}
    \includegraphics[width=4cm]{figures/uoc}
    \quad
    \includegraphics[width=3cm]{figures/uwa}
  \end{center}
}

\usepackage{amsmath}
\usepackage{amssymb}
\usepackage{graphicx}
\usepackage{hyperref}

\usepackage{caption}
\usepackage{subcaption}

% no figure prefixes
\captionsetup[figure]{labelformat=empty,labelsep=none}
\captionsetup[subfigure]{labelformat=empty,labelsep=none}

% matplotlib colors
\definecolor{tab_blue}{HTML}{1F77B4}
\definecolor{tab_orange}{HTML}{FF7F0E}
\definecolor{tab_green}{HTML}{2CA02C}
\definecolor{tab_gray}{HTML}{7F7F7F}

\usepackage[maxnames=6]{biblatex}
\addbibresource{{bibliography.bib}}

\usepackage{tikz}

\newcommand{\dee}{\mathrm{d}}
\newcommand{\E}{\mathsf{E}}
\newcommand{\bm}{\boldsymbol}
\newcommand{\mb}{\mathbf}
\newcommand{\given}{\:\vert\:}

\title{Statistical Finite Elements for Nonlinear PDEs}
\author{Connor Duffin}
\institute[both]{
  Research Assistant\\
  Department of Engineering\\
  University of Cambridge\\
  {\tiny \textit{and}} \\
  PhD Student\\
  Department of Mathematics and Statistics\\
  University of Western Australia

}

\date{MMLDT 2021, 28 September 2021}

\begin{document}

\frame{\titlepage}

\begin{frame}
  \frametitle{Coauthors/reference}

  \begin{figure}[t]
    \centering
    \begin{subfigure}[b]{0.2\linewidth}
      \centering
      \includegraphics[width=\linewidth]{images/ed.jpg}
      \caption{\centering Ed Cripps\\ (UWA)}
    \end{subfigure}
    ~~~
    \begin{subfigure}[b]{0.215\linewidth}
      \centering
      \includegraphics[width=\linewidth]{images/stemler}
      \caption{\centering Thomas Stemler\\ (UWA)}
    \end{subfigure}
    ~~~
    \begin{subfigure}[b]{0.2\linewidth}
      \centering
      \includegraphics[width=\linewidth]{images/mark}
      \caption{\centering Mark Girolami\\ (Cambridge)}
    \end{subfigure}
  \end{figure}
  
  \begin{center}
    \fullcite{duffinStatisticalFiniteElements2021}
  \end{center}
\end{frame}

\begin{frame}
  \frametitle{Background}
  In this work we are interested an internal wave model \alert{(KdV equation)}
  \[
    \partial_t u + \alpha u u_x + \beta u_{xxx} + c u_x = 0
  \]
  with observed data $y_{1:n} = (y_1, \ldots, y_n)$, with PDE parameters $\Lambda = (\alpha,
  \beta, c)$.

  \begin{itemize}
  \item How to use data to \alert{correct for model mismatch?}
  \item Inverse problem: $p(\Lambda \given y_{1:n})$, where the PDE solution $u
    := u(\Lambda)$, \alert{well-studied}.
  \item What if, instead, we look at $p(u^n \given y_{1:n}, \Lambda)$
  \item \alert{How to construct this?}
  \end{itemize}
  \pause
  The \alert{statistical finite element method} (statFEM) is one way.
\end{frame}

\begin{frame}
  \frametitle{Underlying dynamical process}
  \begin{itemize}
  \item Use \alert{stochastic dynamics}. Take the KdV equation:
    \[
      \partial_t u + \alpha u u_x + \beta u_{xxx} + c u_x = \xi_\theta,
    \]
    we have $u := u(x, t)$, $x \in \Omega$, $t \in [0, T]$.
  \item Model \alert{uncertainty} with a Gaussian process,
    \[
      \xi_\theta(x, t) \sim \mathcal{GP}(0, \delta(t - t') \cdot k_\theta(x, x')),
      \quad k_\theta(x, x') = \rho^2 \exp(-\lVert x - x' \rVert^2 / (2\ell^2)).
    \]
  \end{itemize}
  \begin{columns}
    \column{0.65\textwidth} 
    In practise we need to \alert{discretize}: we use \alert{FEM.}
    \begin{itemize}
    \item Construct a mesh, $\Omega_h$, with mesh-size $h$: \alert{discrete
        approximation} to $\Omega$.
    \item Use finite elements on the mesh $\Omega_h$, $u(x, n
      \Delta_t)$ $\approx \sum_{i = 1}^{n_u} u_i \phi_i(x)$.
    \item \alert{Hat-functions} $\implies \phi_i(x_j) = \delta_{ij}$.
    \item \alert{FEM coefficients} $(u_1(t), \ldots, u_{n_u}(t)) \in \mathbb{R}^{n_u}$.
    \end{itemize}
    \column{0.3\textwidth}
    \begin{figure}[t]
      \centering
      \includegraphics[width=\textwidth]{figures/fem-illustrated}
      \caption{FEM \footnotesize{(from Bakka, H. (2019),
          arXiv:1803.03765)}.}
    \end{figure}
  \end{columns}
\end{frame}


\begin{frame}
  \frametitle{Underlying dynamical process}
  \begin{itemize}
  \item Multiply by testing function $\phi_j$, and integrate over $\Omega_h$:
    \[
      \langle \partial_t u_h, \phi_j \rangle
      + \langle \alpha u_h u_{h, x}, \phi_j \rangle
      - \langle \beta u_{h, xx}, \phi_j \rangle
      + \langle c u_{h, x}, \phi_j \rangle = \langle \xi_\theta, \phi_j \rangle, \; j
      = 1, \ldots, n_u.
    \]
    \alert{Removes} higher-order differentiability, so we have \alert{weak
      solutions.} Note: $\langle f, g \rangle = \int_\Omega f \, g \, \mathrm{d} x$.
  \item Time discretization: use \alert{Crank-Nicolson} for stability. Write
    $u^n = (u_1(n\Delta_t), \ldots, u_{n_u}(\Delta_t)) \in \mathbb{R}^{n_u}$.
  \item For the \alert{FEM coefficients} we have a Markov process:
    \[
      \mathcal{M}(u^n, u^{n - 1}) = e_\theta^n, \quad e_\theta^n \sim
      \mathcal{N}(0, \Delta_t G_\theta),
    \]
    for $G_{\theta, ij} = \langle \phi_i, \langle k_\theta(\cdot, \cdot), \phi_j \rangle \rangle$.
  \end{itemize}
  \vspace{0.5cm}
  \pause
  
  How to coherently incorporate data with this (implicit) prior distribution? 
\end{frame}


\begin{frame}
  \begin{block}{{\Large Model-data synthesis}}
    With data $y_n \in \mathbb{R}^{n_y}$, we write $y_n = H u^n + \eta_n$.
    \begin{itemize}
    \item \alert{Observation operator} $H : \mathbb{R}^{n_u} \to \mathbb{R}^{n_y}$.
    \item \alert{Noisy measurements} $\eta_n \sim \mathcal{N}(0, \sigma_n^2 I)$.
    \end{itemize}
    And this gives the state-space model
    \begin{align*}
      \mathcal{M}(u^n, u^{n - 1}) = e_\theta^n ,
      \quad e_\theta^n \sim \mathcal{N}(0, \Delta_t G_\theta),
      \quad &\text{(Dynamics)} \\
      y_n = H u^n + \eta_n ,
      \quad \eta_n \sim \mathcal{N}(0, \sigma_n^2 I).
      \quad &\text{(Data)}
    \end{align*}
    Compute the posterior $p(u^n \given y_{1:n}, \Lambda, \theta)$ using
    nonlinear filtering methods: \alert{ensemble Kalman filter \footfullcite{evensenEnsembleKalmanFilter2003}.}
  \end{block}
\end{frame}


\begin{frame}
  \frametitle{Case study: internal waves in a tub}
  Apply statFEM methodology to measurements of internal waves $u^n$ at $3$
  locations, observing data each of the $n_t = 1001$ timesteps for $t \in [0, 300]$ s.

  \begin{columns}
    \column{0.4\linewidth}
    \begin{figure}[t]
      \centering \newcommand{\twidth}{6cm} \newcommand{\theight}{1.5cm}
      \begin{tikzpicture}
        \coordinate [] (A) at (0, 0); \coordinate [] (B) at (0, \theight);
        \coordinate [] (C) at (-\twidth, 0); \coordinate [] (D) at (-\twidth,
        \theight);

        \coordinate [label=WG3] (z) at (-0.25 * \twidth, \theight); \coordinate
        [label=WG2] (y) at (-0.5 * \twidth, \theight); \coordinate [label=WG1]
        (x) at (-0.75 * \twidth, \theight);

        \draw [gray, dashed] (-0.85 * \twidth, 0.25 * \theight) node[black,
        right] {\footnotesize{$\vartheta^\circ$}} -- (-\twidth, 0.25 *
        \theight); \draw [gray, very thick] (-\twidth, 0.25 * \theight) -- (0,
        0.75 * \theight); \draw [gray, dashed] (x) -- (-0.75 * \twidth, 0);
        \draw [gray, dashed] (y) -- (-0.50 * \twidth, 0); \draw [gray, dashed]
        (z) -- (-0.25 * \twidth, 0); \draw [very thick] (A) -- (C) -- (D) -- (B)
        -- (A);

        \draw [<->] (-\twidth, -0.25) -- node[below] {$L$} (0, -0.25); \draw
        [<->] (-6.25, 0) -- node[left] {$H$} (-6.25, \theight);

        \draw [<->] (-0.52 * \twidth, 0) -- node[left] {$h_2$} (-0.52 * \twidth,
        0.50 * \theight); \draw [<->] (-0.52 * \twidth, 0.50 * \theight) --
        node[left] {$h_1$} (-0.52 * \twidth, \theight);
      \end{tikzpicture}
      \caption{ Wave gauges WG1, WG2, WG3 capture observations of internal wave
        amplitude. Initial conditions are shown as a grey line, labelled with
        initial angle $\vartheta^\circ$. }
    \end{figure}

    \column{0.45\linewidth}
    \begin{figure}[t]
      \centering \includegraphics[width=\textwidth]{figures/raw-processed-data}
      \caption{Experimental data at WG1, WG2, WG3, up to time $t = 300$s.}
      \label{fig:expdata-setup}
    \end{figure}
  \end{columns}
\end{frame}


\begin{frame}
  \frametitle{Results}
  \begin{figure}[t]
    \centering \includegraphics[width=0.8\textwidth]{figures/expdata-posterior}
    \caption{StatFEM posterior measure
      $p(u^n \given y_{1:n}, \Lambda, \theta)$ for the KdV equation.
      (A): posterior at WG locations; (B): posterior wave profile $u$ for
      $t = \{75, 150, 225\}$ s.}
  \end{figure}
\end{frame}


\begin{frame}
  \frametitle{References}
  \printbibliography[heading=none]
\end{frame}

\end{document}
